\subsection{Methodology}

Moving away from traditional Scrum ceremonies to brief meetings every two days has been a beneficial decision for the project. This approach enables continuous project monitoring and discards traditional ceremonies that provide low value in a single-developer context. 

Moving away from traditional Scrum ceremonies to brief meetings every two days has been a beneficial decision for the project. This approach enables continuous project monitoring and eliminates traditional ceremonies that provide low value in a single-developer context.

Iterative meetings with the Product Owner (PO) allowed the project progress to stay on track and enabled the client's feedback to be received at an early stage. Possible delays were detected early, and new features were presented to the PO as soon as possible. At the same time, limiting meeting duration to 30 minutes helped focus on critical issues and prevented time overruns.

Regarding the WIP limit, it is not a restrictive constraint that blocks development progress. Quite the opposite, it has been a useful approach. Limiting the number of active tasks reduced the scope of development, allowing for a more dynamic project flow.

The adoption of ScrumBan for this project has been key to its successful execution so far. By combining selected characteristics of both Scrum and Kanban, this approach is well-suited to the project's needs. Thanks to this methodology, the project's actual execution has closely adhered to the scheduled timeline.
